\documentclass[11pt]{article}
\usepackage[T1]{fontenc}
\usepackage[margin=1in]{geometry}
\usepackage{enumitem}
\usepackage{hyperref}
\setlength{\parindent}{1.5em}
\setlength{\parskip}{6pt}
% =====================================================================
%  EDIT YOUR DETAILS HERE — this ONE file feeds every abstract & deck.
%  (Change a value, recompile; all 36 documents update.)
% =====================================================================
\newcommand{\seminarName}{Aflah Muhammed P}      % <-- your full name (guessed from your email; please verify)
\newcommand{\seminarRoll}{TVE00CS000}            % <-- your university register/roll number
\newcommand{\seminarClassRoll}{00}               % <-- your class roll no (the short one)
\newcommand{\seminarGuide}{Prof.\ <Guide Name>}  % <-- your guide
\newcommand{\seminarDept}{Department of Computer Science and Engineering}
\newcommand{\seminarCollege}{College of Engineering Trivandrum}
\newcommand{\seminarShortCollege}{CET}           % shown in the slide footer, e.g. "Aflah Muhammed P (CET)"
\newcommand{\seminarShortTitle}{B.Tech Seminar}  % middle of the slide footer
\newcommand{\seminarDate}{July 31, 2026}
% ---------------------------------------------------------------------
% Optional: put a logo image named  cet_logo.png  in this folder and it
% will appear on every title slide automatically. If absent, it is skipped.
% =====================================================================

\begin{document}
\begin{center}
{\LARGE\bfseries Multi-Agent Collaboration Mechanisms in LLMs: A Survey}\\[1.1em]
{\large\bfseries Seminar Abstract}\\[0.6em]
{\normalsize \seminarDate}
\end{center}
\vspace{0.6em}
\section*{Abstract}
Large language models have become capable individual agents, yet many real-world tasks -- software engineering, scientific reasoning, negotiation, and societal simulation -- exceed what a single model can plan, remember, and execute alone. A lone LLM is constrained by limited context, hallucination, and a lack of role specialisation, and naively scaling one model yields diminishing returns. Multi-agent systems (MAS) composed of LLM-based agents promise division of labour, debate, and tool use, but the literature is fragmented: individual frameworks such as ChatDev, MetaGPT, AutoGen, and CAMEL each describe bespoke designs without a shared vocabulary for \emph{how} agents actually collaborate, making results difficult to compare, reuse, or generalise.

This survey by Tran et al. (arXiv:2501.06322, 2025) places collaboration at the centre of analysis, shifting attention from isolated models to collaboration-centric design. It formalises an agent as a compact tuple of model, objective, environment, input, and output, and proposes an extensible framework that characterises any collaboration mechanism along five dimensions: the \emph{actors} involved, the \emph{type} of interaction (cooperation, competition, or coopetition), the organisational \emph{structure} (peer-to-peer, centralized, distributed, or hierarchical), the \emph{strategy} (rule-based, role-based, or model-based), and the \emph{coordination protocol}. Using role-based exemplars such as a simulated software company, it shows how these choices shape emergent behaviour, and it reviews applications spanning 5G/6G networks, Industry 5.0, and question answering. This seminar walks through the taxonomy, its illustrative frameworks, and the open challenges and lessons the authors identify on the path toward artificial collective intelligence.
\section*{Key Reference Papers}
\begin{enumerate}
  \item K.-T. Tran, D. Dao, M.-D. Nguyen, Q.-V. Pham, B. O'Sullivan, and H. D. Nguyen, ``Multi-Agent Collaboration Mechanisms: A Survey of LLMs,'' arXiv:2501.06322, 2025.
  \item C. Qian, W. Liu, H. Liu, et al., ``ChatDev: Communicative Agents for Software Development,'' in Proc. ACL, 2024; arXiv:2307.07924.
  \item S. Hong, X. Zheng, J. Chen, et al., ``MetaGPT: Meta Programming for a Multi-Agent Collaborative Framework,'' in Proc. ICLR, 2024; arXiv:2308.00352.
\end{enumerate}
\vspace{0.8em}
\noindent\textbf{Submitted By}\\
\seminarName\\
\seminarRoll

\vspace{0.6em}
\noindent\textbf{Guide}\\
\seminarGuide
\end{document}
